% ============================================================
% Chapter 8: Sustainability and Ethical Implications
%   Section: Ethics of Data Management and Analysis
%     Subsection: Responsibility for Automated Verdicts
% Included from chapters/08-sustainability-ethics/02-data-ethics/data-ethics.tex
% ============================================================

\subsection{Responsibility for Automated Verdicts}

An automated verdict that a claim is false can harm whoever made the
claim, and one that a claim is true can mislead whoever reads it. The
system is built to keep responsibility with people and to make its own
contribution inspectable.

\begin{itemize}
    \item \textbf{Failing towards caution.} When the evidence does not
    decide, the answer is \texttt{UNVERIFIED}; a definitive verdict that
    cites nothing, or only sources the model itself called unrelated, is
    downgraded; a \texttt{TRUE} confirmed by a single outlet keeps a
    capped confidence.
    \item \textbf{Failures are not verdicts.} A claim whose search or
    model failed is marked as such, and the reader view says ``not
    checked'' rather than showing the \texttt{UNVERIFIED} it is stored
    as. A search outage is never presented as a judgement about a claim.
    \item \textbf{Every verdict shows its work.} The model's reasoning,
    its answer before the guardrails changed it, each source's stance and
    quoted passage, and whether each source's reliability is a real
    rating or a default, are visible to the reader.
    \item \textbf{A person decides what is published.} The system
    proposes and the editor chooses. Its accuracy has been measured on
    one benchmark and found low, and it is not presented as a
    fact-checker whose verdicts stand on their own.
\end{itemize}
